\chapter{Results}
\label{chap:results}

\section{Introduction}

This chapter presents the empirical findings from the 2022 Ghana Demographic and Health Survey. It begins with the construction of the analytic sample and weighted descriptive patterns of childhood stunting. It then presents the Bayesian hierarchical modelling results. At the present stage, the descriptive analysis is complete and validated against the published GDHS patterns, while the final three-level Bayesian models are being fitted using the reproducible scripts described in Chapter Three.

\section{Analytic Sample}

The analytic sample contained 4,928 children aged 0--59 months with valid height-for-age measurements who slept in the sampled household on the night preceding the survey, consistent with the DHS de facto anthropometry population. Among these children, 927 were classified as stunted using the WHO criterion of height-for-age z-score below -2 standard deviations. The sum of the rescaled DHS sampling weights for these children was 4,292.97, which rounds to the weighted denominator of 4,293 children reported for height-for-age in Table 11.1 of the official 2022 GDHS. This agreement provides an additional validation of the analytic-sample construction and weighting procedure \citep{GSSICF2024}.

The children were distributed across 616 survey communities and 3,545 households. Of the 3,545 households represented in the analytic sample, 2,398 contained one eligible child, 951 contained two, 165 contained three, 26 contained four, four contained five, and one household contained nine eligible children. In total, 1,147 households contained more than one eligible child. Approximately 51.3\% of children in the analytic sample lived in households containing at least two eligible children. This degree of within-household clustering provides direct empirical justification for including a household-level random effect in the hierarchical model.

The mean number of eligible children per household was 1.39, while the mean number of eligible children per community was 8.0. Community sample sizes ranged from 1 to 28 children.

\section{Overall Prevalence of Childhood Stunting}

The unweighted prevalence of stunting in the analytic sample was 18.8\%. After application of the DHS household-member sampling weights, the estimated prevalence was 17.4\% (95\% CI: 15.8\%--18.9\%). The Taylor-linearized standard error was 0.008 after rounding, matching the value reported in Appendix B of the 2022 GDHS. This agreement provides an important validation of the analytic sample construction, weighting, and survey-design specification \citep{GSSICF2024}.


\begin{table}[H]
\centering
\caption{Weighted prevalence of childhood stunting by selected characteristics}
\label{tab:descriptive_stunting}
\begin{tabular}{lrrl}
\toprule
Characteristic & Unweighted $n$ & Stunted $n$ & Weighted prevalence (95\% CI) \\
\midrule
Overall & 4,928 & 927 & 17.4 (15.8--18.9) \\
\addlinespace
Male & 2,529 & 529 & 19.4 (17.2--21.6) \\
Female & 2,399 & 398 & 15.3 (13.3--17.3) \\
\addlinespace
Urban & 2,024 & 337 & 15.1 (12.6--17.5) \\
Rural & 2,904 & 590 & 19.6 (17.6--21.6) \\
\addlinespace
Poorest & 1,601 & 375 & 22.8 (19.7--25.9) \\
Poorer & 1,220 & 275 & 22.0 (19.1--24.9) \\
Middle & 877 & 134 & 16.9 (13.3--20.5) \\
Richer & 694 & 103 & 15.5 (11.7--19.3) \\
Richest & 536 & 40 & 7.5 (4.7--10.3) \\
\bottomrule
\end{tabular}

\vspace{0.4em}
\begin{minipage}{0.93\textwidth}
\footnotesize
\textit{Note:} Prevalence estimates use DHS sampling weights. Confidence intervals use Taylor linearization with survey strata and primary sampling units.
\end{minipage}
\end{table}

\section{Stunting by Age of the Child}

Weighted stunting prevalence varied substantially by age. Prevalence was 13.2\% among children aged 0--5 months and 11.5\% among those aged 6--11 months. It increased to 17.2\% among children aged 12--23 months and reached its highest level, 24.6\%, among children aged 24--35 months. Prevalence then declined to 18.8\% among children aged 36--47 months and 14.2\% among those aged 48--59 months.

The descriptive age pattern is consistent with accumulated linear growth faltering during the second and third years of life, with the highest observed burden occurring among children aged 24--35 months.


\begin{figure}[H]
\centering
\includesvg[width=0.88\textwidth]{stunting_by_age}
\caption{Survey-weighted stunting prevalence by child age group.}
\label{fig:stunting_age}
\end{figure}

\section{Stunting by Sex}

Weighted stunting prevalence was higher among male than female children. The estimated prevalence was 19.4\% among boys compared with 15.3\% among girls. This difference is consistent with the direction observed in several previous studies of childhood stunting in sub-Saharan Africa and Ghana.

\section{Stunting by Household Wealth}

A pronounced socioeconomic gradient was observed across household wealth quintiles. Weighted stunting prevalence was 22.8\% among children in the poorest households and 22.0\% among those in the poorer quintile. Prevalence declined to 16.9\% in the middle quintile, 15.5\% in the richer quintile, and 7.5\% among children in the richest households.

The descriptive results therefore indicate a strong inverse relationship between household wealth and stunting prevalence, although the multivariable hierarchical models are required to determine the extent to which this gradient persists after adjustment for other characteristics.


\begin{figure}[H]
\centering
\includesvg[width=0.88\textwidth]{stunting_by_wealth}
\caption{Survey-weighted stunting prevalence by household wealth quintile.}
\label{fig:stunting_wealth}
\end{figure}

\section{Stunting by Place of Residence}

Children living in rural areas had a weighted stunting prevalence of 19.6\%, compared with 15.1\% among children living in urban areas. The difference suggests a rural disadvantage at the descriptive level. However, rural residence is strongly associated with household wealth, region, access to water and sanitation, and other contextual conditions, and should therefore not be interpreted independently of the multivariable analysis.

\section{Regional Variation}

Substantial regional differences were observed. The highest weighted prevalence occurred in the Northern Region at 29.6\% and the North East Region at 29.3\%. Stunting prevalence was also relatively high in the Upper East Region and Savannah Region, both at approximately 21\%, and in Oti Region at 20.3\%.

The lowest weighted prevalence estimates were observed in Eastern Region at 10.4\%, Western North Region at 10.6\%, and Greater Accra Region at 11.4\%. These descriptive differences closely resemble the regional pattern reported in the official 2022 GDHS and reinforce the importance of accounting for contextual heterogeneity.


\begin{figure}[H]
\centering
\includesvg[width=0.88\textwidth]{stunting_by_region}
\caption{Survey-weighted stunting prevalence by region.}
\label{fig:stunting_region}
\end{figure}

\section{Maternal Education}

Among children for whom maternal education information was available, weighted stunting prevalence decreased markedly with increasing maternal educational attainment. Prevalence was 24.4\% among children whose mothers had no formal education, 18.5\% among those whose mothers had primary education, 15.3\% among those whose mothers had secondary education, and 5.7\% among those whose mothers had higher education.

Maternal education was unavailable for approximately 8.6\% of children in the preliminary analytic sample. For this reason, maternal education will be introduced in a later model using a harmonized sample so that changes in coefficients are not confounded by changes in the observations analysed.


\begin{figure}[H]
\centering
\includesvg[width=0.88\textwidth]{stunting_by_maternal_education}
\caption{Survey-weighted stunting prevalence by maternal educational attainment. The missing category represents children without linked maternal-education information.}
\label{fig:stunting_matedu}
\end{figure}

\section{Water and Sanitation Patterns}

Water and sanitation were summarized using JMP-consistent improved/unimproved facility-type groupings. Weighted stunting prevalence was 16.4\% among children in households using an improved drinking-water source and 23.0\% among those using an unimproved source. For sanitation, weighted stunting prevalence was 13.8\% among children in households with an improved sanitation facility type and 22.3\% among those with an unimproved facility or no facility.

These classifications refer to source or facility type only. They do not represent the full JMP basic or safely managed service ladders because the latter require additional information such as collection time, sharing, availability, management, or water quality.

These crude differences should not be interpreted causally because WASH conditions are strongly correlated with wealth, residence, and region.

\begin{figure}[H]
\centering
\includesvg[width=0.78\textwidth]{stunting_by_sanitation}
\caption{Survey-weighted stunting prevalence by improved versus unimproved/no sanitation facility type.}
\label{fig:stunting_sanitation}
\end{figure}


\begin{figure}[H]
\centering
\includesvg[width=0.88\textwidth]{stunting_by_water_source}
\caption{Survey-weighted stunting prevalence by improved versus unimproved drinking-water source type.}
\label{fig:stunting_water}
\end{figure}

\section{Household and Community Clustering}

The structure of the analytic data demonstrates that stunting observations are not statistically independent. More than half of the children belonged to households containing at least one additional eligible child, and all children were nested within survey communities. An independent GEE diagnostic also indicated substantially stronger residual dependence within households than within communities, with exchangeable working correlations of approximately 0.154 and 0.024, respectively. These GEE quantities are not interchangeable with the Bayesian variance components reported below; they are used only as an independent check on the presence and relative scale of clustering.

\section{Bayesian Hierarchical Model Results}

Three Bayesian hierarchical extensions were examined after the community-only pilot. Model 1 included both household and community random intercepts alongside the core child, household-wealth, residence, and regional covariates. Model 2 added drinking-water and sanitation source/facility type. Model 3 added maternal education and was fitted to the 4,503 children with linked maternal-education information. To separate the effect of maternal education from the effect of changing the analysis sample, Model 2 was also refitted on the same 4,503-child harmonized sample.

Across all of these fits, no divergent NUTS transitions were observed. Fixed-effect coefficients showed good cross-chain agreement, with $\hat R$ values close to 1.00. The hierarchical standard-deviation parameters mixed more slowly than the fixed effects; maximum $\hat R$ values were approximately 1.018 for Model 1, 1.012 for Model 2 on the full sample, 1.036 for the harmonized Model 2, and 1.033 for Model 3. Consequently, the direction and magnitude of the fixed-effect findings are considered stable, while the exact higher-level variance estimates will be confirmed with longer production runs before final thesis submission.

\subsection{Model 1: Household and Community Random Effects}

Introducing a household random intercept revealed substantially more residual heterogeneity between households than between communities. The posterior median community-level standard deviation was 0.416 (95\% posterior interval: 0.159--0.590), whereas the household-level standard deviation was 1.066 (0.774--1.346).

On the latent logistic scale, the median community ICC was 0.038, while the household-specific variance partition coefficient was 0.246. The correlation for two otherwise similar children from the same household, which includes both shared household and shared community variation, was 0.285. Expressed on the odds-ratio scale, the community MOR was 1.49 and the household MOR was 2.76. Thus, after adjustment for the measured covariates, unexplained differences between households were substantially larger than unexplained differences between communities.

The fixed-effect pattern was broadly consistent with the earlier pilot. Boys had higher posterior odds of stunting than girls (OR 1.45, 95\% posterior interval: 1.23--1.72). Children aged 24--35 months had approximately three times the odds of stunting compared with children aged 0--5 months (OR 2.95, 2.10--4.20). Relative to the poorest quintile, the middle (OR 0.58, 0.42--0.79), richer (OR 0.54, 0.38--0.78), and richest (OR 0.25, 0.15--0.38) groups had lower posterior odds. Rural residence was not clearly separated from urban residence after adjustment (OR 0.89, 0.69--1.13).

\subsection{Model 2: Addition of Water and Sanitation}

In the full-sample WASH model, children living in households using an unimproved drinking-water source had higher posterior odds of stunting than those using an improved source (OR 1.36, 95\% posterior interval: 1.05--1.78). The posterior probability that the corresponding log-odds coefficient was positive was approximately 0.989. In contrast, the sanitation comparison was less decisive after adjustment: the OR for an unimproved or absent sanitation facility was 1.16 (0.91--1.45), with posterior probability of a positive coefficient of approximately 0.872.

Adding the WASH variables did not eliminate the hierarchical structure. The posterior median community SD changed from 0.416 in Model 1 to 0.390 in Model 2, while the household SD changed from 1.066 to 1.100. The corresponding household VPC remained substantial at approximately 0.260. These changes are small relative to the posterior uncertainty and do not support the interpretation that the WASH variables explain most of the residual household heterogeneity.

\subsection{Model 3: Addition of Maternal Education}

For the maternal-education analysis, Model 2 was first refitted to the 4,503 children with non-missing maternal education. On this harmonized sample, unimproved water remained associated with higher posterior odds of stunting (OR 1.51, 1.16--1.98), whereas the sanitation comparison remained less clearly separated from one (OR 1.13, 0.87--1.45).

After maternal education was added in Model 3, the posterior association for unimproved drinking water remained essentially unchanged (OR 1.50, 1.14--2.00), with $P(\text{OR}>1\mid\text{data})\approx0.998$. The sanitation estimate remained comparatively uncertain (OR 1.11, 0.87--1.43; $P(\text{OR}>1\mid\text{data})\approx0.79$).

Maternal higher education showed the clearest education association. Relative to children whose mothers had no formal education, children whose mothers had higher education had substantially lower posterior odds of stunting (OR 0.38, 0.19--0.70), with $P(\text{OR}<1\mid\text{data})\approx0.999$. Primary education was not clearly separated from no education (OR 1.04, 0.77--1.42), and the secondary-education estimate was also compatible with a range of modest protective or null associations (OR 0.90, 0.70--1.17).

The wealth gradient attenuated after maternal education entered the model. For example, the richest-versus-poorest OR changed from 0.26 in the harmonized WASH model to 0.36 in Model 3, although the posterior interval remained clearly below one (0.20--0.63). This attenuation is consistent with overlap between maternal education and household socioeconomic position, but the cross-sectional design does not permit a causal mediation interpretation.

\begin{table}[H]
\centering
\caption{Selected posterior odds ratios from the hierarchical models}
\label{tab:key_posterior_or}
\begin{tabular}{lll}
\toprule
Model and comparison & Posterior median OR & 95\% posterior interval \\
\midrule
Model 1: Male vs female & 1.45 & 1.23--1.72 \\
Model 1: Age 24--35 vs 0--5 months & 2.95 & 2.10--4.20 \\
Model 1: Richest vs poorest & 0.25 & 0.15--0.38 \\
Model 2 full: Unimproved water & 1.36 & 1.05--1.78 \\
Model 2 full: Unimproved/no sanitation & 1.16 & 0.91--1.45 \\
Model 3: Unimproved water & 1.50 & 1.14--2.00 \\
Model 3: Unimproved/no sanitation & 1.11 & 0.87--1.43 \\
Model 3: Maternal primary vs none & 1.04 & 0.77--1.42 \\
Model 3: Maternal secondary vs none & 0.90 & 0.70--1.17 \\
Model 3: Maternal higher vs none & 0.38 & 0.19--0.70 \\
\bottomrule
\end{tabular}
\end{table}

\begin{figure}[H]
\centering
\includesvg[width=0.95\textwidth]{model3_posterior_or_forest}
\caption{Selected posterior odds ratios from Model 3 with 95\% posterior intervals.}
\label{fig:model3_forest}
\end{figure}

\subsection{Comparison of Household and Community Heterogeneity Across Models}

Table~\ref{tab:variance_comparison} summarizes the random-effect results. Household heterogeneity remained substantially larger than community heterogeneity throughout the model sequence. On the full analytic sample, the community SD decreased slightly after WASH was added, from 0.416 to 0.390, whereas the household SD remained close to or above 1.0. On the harmonized maternal-education sample, the household SD was approximately 1.22 before adding maternal education and 1.24 after it was added. Given the overlapping posterior intervals and slower mixing of these variance parameters, the small numerical increases should not be interpreted as evidence that maternal education increases unexplained heterogeneity. Rather, the central finding is that substantial household-level residual variation persists after WASH and maternal education are included.

\begin{table}[H]
\centering
\caption{Posterior summaries of hierarchical heterogeneity across models}
\label{tab:variance_comparison}
\begin{tabular}{lcccc}
\toprule
Model & Community SD & Household SD & Community ICC & Household VPC \\
\midrule
Model 1 & 0.416 & 1.066 & 0.038 & 0.246 \\
Model 2 full & 0.390 & 1.100 & 0.033 & 0.260 \\
Model 2 harmonized & 0.391 & 1.222 & 0.031 & 0.301 \\
Model 3 & 0.411 & 1.239 & 0.034 & 0.307 \\
\bottomrule
\end{tabular}

\vspace{0.4em}
\begin{minipage}{0.94\textwidth}
\footnotesize
\textit{Note:} SD denotes the latent-scale random-intercept standard deviation. ICC is the community variance share on the latent logistic scale. Household VPC is the household-specific latent variance share. Model 2 harmonized and Model 3 use the same 4,503-child analysis sample.
\end{minipage}
\end{table}

\begin{figure}[H]
\centering
\includesvg[width=0.95\textwidth]{model_variance_comparison}
\caption{Posterior medians and 95\% posterior intervals for household and community random-effect standard deviations across models.}
\label{fig:variance_comparison}
\end{figure}

\section{Current Diagnostic Status and Remaining Sensitivity Analyses}

The fixed-effect parameters in Models 1--3 showed satisfactory chain agreement and no divergent transitions were observed. The hierarchical variance parameters were more computationally demanding, particularly in the harmonized Model 2 and Model 3, where $\hat R$ for the community SD was approximately 1.036 and 1.033, respectively. These values do not overturn the broad conclusion that household heterogeneity is substantially larger than community heterogeneity, because the household/community contrast is large and consistent across models, but they indicate that longer production runs are warranted before final submission.

The results in this section therefore represent the current inferential analysis rather than the final locked publication analysis. Prior-sensitivity analysis, survey-weight sensitivity, expanded posterior predictive checks, predictive model comparison, and longer-chain confirmation of the hierarchical variance components remain to be completed before the final discussion and conclusions are written.

\section{Chapter Summary}

The weighted descriptive analysis estimated a national childhood stunting prevalence of 17.4\%. The hierarchical models showed that residual heterogeneity is concentrated more strongly at the household level than at the community level. In Model 1, the household VPC was approximately 0.25 compared with a community ICC of approximately 0.04, and the household MOR was about 2.76 compared with a community MOR of about 1.49.

After adjustment for demographic, socioeconomic, regional, household, and community structure, unimproved drinking-water source remained associated with higher posterior odds of stunting. Sanitation showed a weaker and less certain adjusted association. Maternal higher education was associated with substantially lower posterior odds of stunting, while primary and secondary education were not clearly separated from no formal education after full adjustment. The household variance remained substantial after WASH and maternal education were introduced, indicating that important unmeasured household-level influences remain. These findings will be subjected to the planned prior, survey-weight, and predictive sensitivity analyses before the final thesis discussion is completed.
